\subsection{Matching a moving soft boundary}
\label{sec:compton-soft-matching-check}

The separate rational subtraction terms can have stronger soft poles than
 their sum.  Their finite parts therefore require the exact phase-space
boundary and the first subleading term of the regulator.  This section derives
that matching directly from the integration region.

Let $H>0$, $s=H+1$, and $y=H-x+z$.  The dimensionless Dalitz region is
\begin{equation}
 0<x<H,\qquad 0<z<\frac{x(H-x)}{s-x}.
\end{equation}
At either soft corner introduce a variable $e$, which here denotes a
kinematic variable rather than the electric charge, and $r=z/e$.  The two
choices are $e=x$ and $e=y$.  In both cases the Jacobian is $e$ and
\begin{align}
 xy&=e[H+e(r-1)],\qquad
 0<e<H,\qquad 0<r<R(e)=\frac{H-e}{s-e},\\
 R(e)&=q+q_1e+O(e^2),\qquad
 q=\frac{H}{s},\qquad q_1=-\frac{1}{s^2}.
\end{align}
Consider a kernel whose expansion at this corner is
\begin{equation}
 K(e,er)=\frac{a(r)}{e^3}+\frac{b(r)}{e^2}+O(e^{-1}),
 \label{eq:compton-matching-kernel}
\end{equation}
where $a$ and $b$ are smooth on $[0,q]$.  Terms with $a=0$ include the usual
logarithmic soft singularity.  Define the four moments
\begin{align}
 P&=\int_0^q r a(r)\,\mathrm{d}r,\\
 A&=\int_0^q b(r)\,\mathrm{d}r+q_1a(q),\\
 B&=\int_0^q b(r)\ln r\,\mathrm{d}r+q_1a(q)\ln q,\\
 C&=\frac{1}{H}\int_0^q(1-r)a(r)\,\mathrm{d}r.
\end{align}
The moment $A$ includes the displacement of the upper boundary.  Indeed,
after multiplication by the Jacobian, the singular part can be written,
in the sense of integration against smooth functions of $r$, as
\begin{equation}
 eK\,\theta(R(e)-r)
 =\frac{a(r)\theta(q-r)}{e^2}
 +\frac{b(r)\theta(q-r)+q_1a(q)\delta(r-q)}{e}
 +O(1).
 \label{eq:compton-moving-boundary-distribution}
\end{equation}
Here $\theta$ is the Heaviside step function and $\delta$ is the
Dirac delta distribution. The last term includes distributions at $r=q$ with integrable coefficients
in $e$.  It contributes the same finite limit in the two regularizations.

Let $I_z(\zeta)$ be the integral with $z>\zeta>0$.  Let $I_\eta$ be the
integral multiplied by $(xy)^\eta$, initially at sufficiently large
$\operatorname{Re}\eta$ and then continued to zero.  Define their finite
parts by
\begin{align}
 I_z(\zeta)&=\frac{P}{\zeta}-A\ln\zeta+F_z+o(1),\\
 I_\eta&=\frac{A}{\eta}+F_\eta+O(\eta).
\end{align}
The relation between them is
\begin{equation}
 \boxed{F_\eta=F_z+A\ln H-B-C.}
 \label{eq:compton-general-soft-matching}
\end{equation}
Contributions from distinct corners are added.  All quantities here refer
to the unnormalized integral, before the common cross-section prefactor.

To prove the relation, first replace the regulator by $(He)^\eta$.
Integration of the $e^{-1}$ term in
\eqref{eq:compton-moving-boundary-distribution} gives $A/\eta+A\ln H$.
Changing the lower boundary from $e>\zeta/r$ to an $r$-independent lower
boundary gives the logarithmic moment $B$.  The $e^{-2}$ term gives
$P/\zeta$, with no additional finite shift for a fixed upper boundary.
The remaining regulator dependence follows from the exact identity for
$xy$:
\begin{equation}
 (xy)^\eta=(He)^\eta
 \left[1-\frac{\eta e(1-r)}{H}
 +O(\eta e^2)+O(\eta^2e^2)\right].
\end{equation}
Its linear term multiplies $a/e^2$.  The resulting power-regulated integral contains
$\eta\int e^{\eta-1}\mathrm{d}e$, whose limit is one.  This gives $-C$.
Terms quadratic in $e$ are integrable at $\eta=0$ and give no finite
regulator shift.  Omitting the linear term would thus give an incorrect
constant even if the residue $A$ were correct.

The corner of the moving boundary does not add another constant.
The lower intersection with $z=\zeta$ satisfies
\begin{equation}
 e_{\min}=\frac{\zeta}{q}-\frac{q_1\zeta^2}{q^3}
 +O(\zeta^3).
\end{equation}
The interval displaced from $\zeta/q$ has width $O(\zeta^2)$.
After the leading angular interval has closed, its integrated density
is at most $O(\zeta^{-1})$, so its contribution is $O(\zeta)$.
Smoothness at $r=0$ likewise makes the excluded strip contribute a
vanishing remainder after the displayed singular terms have been removed.
This establishes that the boundary term already present in $A$ and $B$
is sufficient.

\paragraph{A logarithmic test kernel.}
For $K=x^{-2}$, $a=0$, $b=1$, and $A=q$.  The exact analytically regulated integral is
\begin{equation}
 I_\eta^{(2)}=
 H^{2\eta}\frac{\Gamma(\eta-1)\Gamma(\eta+1)}{\Gamma(2\eta+1)}
 \left[{}_2F_1(\eta+1,\eta-1;2\eta+1;q)-1\right].
\end{equation}
This follows by first integrating $y^\eta$ between
$H-x$ and $s(H-x)/(s-x)$ and then using the Euler beta integral.
Let
\begin{equation}
 x_\pm=\frac{H+\zeta\pm\sqrt{(H-\zeta)^2-4\zeta}}{2}.
\end{equation}
The independent cutoff integral and the two finite parts are
\begin{align}
 I_z^{(2)}&=q\ln\frac{x_+}{x_-}
 +\frac{1}{s}\ln\frac{s-x_+}{s-x_-}
 +\zeta\left(\frac{1}{x_+}-\frac{1}{x_-}\right),\\
 F_z^{(2)}&=2q\ln H-\ln s-q,\\
 F_\eta^{(2)}&=2q\ln H-\frac{\ln s}{s}.
\end{align}
Since $B=q(\ln q-1)$, these expressions obey
\eqref{eq:compton-general-soft-matching}.  In particular, the finite
$-q$ in the cutoff integral cannot be discarded.

\paragraph{A power-divergent test kernel.}
For $K=x^{-3}$, $a=1$, $b=0$, and the residue comes entirely from the
moving boundary: $A=-1/s^2$.  Direct integration gives
\begin{align}
 I_\eta^{(3)}&=
 H^{2\eta-1}\frac{\Gamma(\eta-2)\Gamma(\eta+1)}{\Gamma(2\eta)}
 \left[{}_2F_1(\eta+1,\eta-2;2\eta;q)-1\right],\\
 I_z^{(3)}&=J(x_+)-J(x_-),\\
 J(x)&=-\frac{q}{x}-\frac{\ln x}{s^2}
 +\frac{\ln(s-x)}{s^2}+\frac{\zeta}{2x^2}.
\end{align}
Their expansions yield
\begin{align}
 P&=\frac{q^2}{2},\qquad B=-\frac{\ln q}{s^2},\qquad
 C=\frac{q-q^2/2}{H},\\
 F_z^{(3)}&=-\frac{1}{s}-\frac{2\ln H}{s^2},\\
 F_\eta^{(3)}&=-\frac{1}{s}-\frac{2\ln H+\ln s}{s^2}
 -\frac{q-q^2/2}{H}.
\end{align}
Both the boundary term and $C$ are necessary for agreement.

\paragraph{The Bose subtraction correction.}
For the correction $D=R_{VW}[T_{UW}(y,z)-T_{VW}(x,z)]$ constructed from the
rational coefficients and their local subtraction terms, the $x$ corner
is integrable.  At the $y$ corner its leading coefficient is
\begin{equation}
 a(r)=\frac{16(H+1)(Hr+2)(Hr-2H-6)}{3H(r-1)^3}.
\end{equation}
The coefficient $b(r)$ is obtained by differentiating $e^3D$ at $e=0$,
after substituting $y=e$, $z=er$, and $x=H+e(r-1)$.
This is a finite algebraic operation on the exact rational subtraction.
The only poles of these coefficients in $r$ are at $1$ and $-4/H$,
outside $[0,q]$.  Their moments are therefore closed logarithms and
dilogarithms.  For example, every repeated pole follows by differentiation
with respect to its location $u$:
\begin{align}
 \int_0^q\frac{\mathrm{d}r}{(r-u)^n}
 &=\frac{\partial_u^{n-1}}{(n-1)!}\ln(1-q/u),\\
 \int_0^q\frac{\ln r\,\mathrm{d}r}{(r-u)^n}
 &=\frac{\partial_u^{n-1}}{(n-1)!}
 \left[\ln q\ln(1-q/u)+\operatorname{Li}_2(q/u)\right].
\end{align}
The arguments make both expressions real for these pole locations.
\paragraph{Integrating the rational correction.}
Write the exact partial fraction decomposition of $D$ in $x$ as
\begin{equation}
 D(x,z)=\sum_{a,n}\frac{c_{a,n}(z)}{(x-a(z))^n}.
\end{equation}
There are nine terms, with $(a,n)$ equal to
\begin{gather}
 (H+z,1),\ (H+z,2),\ (H,1),\ (H,2),\ (H,3),\
 (z,1),\ (z,2),\ (a_*,1),\ (a_*,2),\\
 a_* = \frac{-4H+z^2-4z}{z-4}.
\end{gather}
The coefficients are fixed algebraically by the residues of $D$, including
 derivatives for repeated poles.  Thus no fit is involved.  Rationalize
the physical boundary by
\begin{equation}
 z=s+1-u-\frac{s}{u},\qquad
 x_-=s-u,\qquad x_+=s-\frac{s}{u},\qquad \sqrt{s}<u<s.
\end{equation}
The inner integral is explicitly
\begin{align}
 \int_{x_-}^{x_+}D\,\mathrm{d}x
 &=\sum_a c_{a,1}(z)\ln\frac{x_+-a}{x_--a}\\
 &\quad+\sum_{a,n>1}\frac{c_{a,n}(z)}{1-n}
 \left[(x_+-a)^{1-n}-(x_--a)^{1-n}\right].
\end{align}
The ratios are positive in the physical interval, as follows directly
from the positions of the poles relative to $[x_-,x_+]$.
After multiplication by $1-s/u^2$, the outer integrand consists of rational
functions multiplying the following three real logarithms:
\begin{align}
 L_1&=\ln\frac{u^2}{s},\\
 L_2&=\ln\frac{u(u-1)}{s-u},\\
 L_3&=\ln\frac{(u-1)(u+s)^2}{u(u+1)^2(s-u)}.
\end{align}
Each logarithm vanishes at $u=\sqrt{s}$.  Partial fractions followed by
integration by parts therefore gives a finite sum of generalized
polylogarithms of weight at most two, based at $\sqrt{s}$, with rational
coefficients.  The letters are the zeros of the rational denominators and
of the logarithm arguments.  Differentiating that finite sum checks the
original outer integrand term by term.

The power divergence also affects the conversion of the outer endpoint.
For $b=s-u$, the exact relation gives
\begin{equation}
 \zeta=\frac{H}{s}b-\frac{b^2}{s^2}+O(b^3),\qquad
 \frac{P}{\zeta}=\frac{Ps}{Hb}+\frac{P}{H^2}+O(b).
\end{equation}
Consequently, if $F_b$ is the finite part after subtracting the
$b^{-1}$ and $\ln b$ terms, then
\begin{equation}
 F_z=F_b+A\ln(H/s)-\frac{P}{H^2}.
\end{equation}
The last constant is required before applying
\eqref{eq:compton-general-soft-matching}.  Convergent Taylor series for the
endpoint remainder and for the generalized polylogarithms evaluate this
closed result without quadrature.
